\documentclass[10pt,a4paper]{amsart}
\usepackage[margin=19mm]{geometry}
\usepackage{amsmath,amssymb,mathtools}
\usepackage[scr=rsfs,cal=euler]{mathalfa}
\usepackage{xfrac}
\usepackage[unicode,hidelinks,bookmarks=false]{hyperref}
\newcommand{\ie}{i.e.\ }
\newcommand{\cf}{cf.\ }
\newcommand{\resp}{respectively\ }
\newcommand{\Subsection}{\subsection}
\providecommand{\nobreakdash}{\nobreak\hbox{-}}
\setlength{\emergencystretch}{2em}
\multlinegap0pt
\usepackage{amssymb}
\usepackage{mathtools}
\usepackage{bm}
\usepackage[shortlabels]{enumitem}
\newtheorem{proposition}{Proposition}
\newtheorem{lemma}{Lemma}
\title{Conformal Parametrisation of Star-Shaped Domains and Persistence of Semilinear Elliptic Solutions}
\author{Gianni Arioli}
\date{Source: arXiv:2609.14546v1 (CC BY 4.0)}
\begin{document}
\maketitle

\noindent\emph{Source and licence.} Conformal Parametrisation of Star-Shaped Domains and Persistence of Semilinear Elliptic Solutions. Version arXiv:2609.14546v1 (CC BY 4.0), distributed by arXiv under the Creative Commons Attribution 4.0 International licence, http://creativecommons.org/licenses/by/4.0/, as recorded in the arXiv licence metadata for this exact version and shown on the abstract page https://arxiv.org/abs/2609.14546v1. Full licence notice: \texttt{licenses/arxiv-2609.14546-CC-BY-4.0.txt}.\par\smallskip

\noindent\textit{Editorial scope.} Counted unit: the article's lemma lem:Cf (source0.tex lines 2975--3001). The article's complete original proof of it is delivered with it (source0.tex lines 3003--3199, one \texttt{proof} environment of 197 lines). No mathematics is written, repaired or re-derived: the statement and the proof are the article's own lines, in the article's own notation, and no retained line is substituted or re-flowed.  Retained uncounted as the source-local context the proof needs: lines 1398--1400; lines 1508--1524.  Nothing was omitted from any retained span, so the package carries no omission record.  No retained line contains a citation, so the wrapper carries no bibliography and no bibliography entry.  No retained line contains a figure environment, an image-inclusion command or drawing code, so no image is delivered and the package declares no asset dependency.  Editorial formatting only, in addition: an AMS \texttt{amsart} preamble that restates the article's own declarations and macros used by the retained lines, namely the environments \texttt{proposition}, \texttt{lemma} on separate counters, together with the standard packages those lines need.  The retained environments print in this document as Proposition 1, Lemma 1 in order of appearance, whereas the article prints the same items with the numbering of its own full document.  The complete native source of the article as distributed in the arXiv e-print for this exact version is bundled unchanged as \texttt{sources/210330/source0.tex}.
\begin{equation}\label{eq:composition-margin}
	\rho e^{\bar d+R}<\tau.
\end{equation}

\begin{proposition}\label{prop:composition-strip}
	Let $g\in\mathcal A_\tau$ and let
	\[
	u=\bar u+h\in\mathcal U.
	\]
	If \eqref{eq:composition-margin} holds, then
	\[
	g(\,\cdot+u(\cdot)\,)\in\mathcal A_\rho
	\]
	and
	\begin{equation}\label{eq:composition-strip}
		\|g(\,\cdot+u(\cdot)\,)\|_\rho
		\le
		C_{\rm comp}\|g\|_\tau.
	\end{equation}
	The estimate is uniform for $u\in\mathcal U$.
\end{proposition}

\begin{lemma}\label{lem:Cf}
	Let $r_1,r_2\in B_\delta(r_0)$, and let $u_1,u_2$ be the
	corresponding solutions of the Theodorsen equation. Assume that
	\[
	\|u_1-u_2\|_\rho
	\le
	C_u\|r_1-r_2\|_\sigma .
	\]
	Then, for every $1<\rho_f<\rho$,
	\begin{equation}\label{eq:f-lipschitz}
		\|f_{r_1}-f_{r_2}\|_{\rho_f}
		\le
		C_f\|r_1-r_2\|_\sigma ,
	\end{equation}
	where
	\begin{equation}\label{eq:Cf}
		C_f
		:=
		\frac{\rho_f}{\rho-\rho_f}\,
		\rho e^d
		\left[
		1+(K_\delta+R_\delta)C_u
		\right],
		\qquad
		d:=\bar d+R .
	\end{equation}
\end{lemma}

\begin{proof}
	Set
	\[
	\Delta_r:=\|r_1-r_2\|_\sigma
	\]
	and, for $z\in S_\rho$, define
	\[
	F_j(z)
	:=
	r_j\bigl(z+u_j(z)\bigr)
	e^{i(z+u_j(z))},
	\qquad j=1,2.
	\]
	On the real axis,
	\[
	F_j(t)=f_{r_j}(e^{it}).
	\]
	Moreover, $F_j$ is holomorphic and $2\pi$-periodic in $S_\rho$.
	Since its boundary values are those of the Riemann map, its
	nonpositive Fourier modes vanish. Hence, by analytic continuation,
	\begin{equation}\label{eq:F-f-identification}
		F_j(z)=f_{r_j}(e^{iz}),
		\qquad z\in S_\rho.
	\end{equation}
	
	We decompose
	\[
	\begin{aligned}
		F_1(z)-F_2(z)
		={}&
		\bigl[
		r_1(z+u_1(z))-r_2(z+u_1(z))
		\bigr]e^{i(z+u_1(z))}
		\\
		&+
		\bigl[
		r_2(z+u_1(z))-r_2(z+u_2(z))
		\bigr]e^{i(z+u_1(z))}
		\\
		&+
		r_2(z+u_2(z))
		\bigl[
		e^{i(z+u_1(z))}-e^{i(z+u_2(z))}
		\bigr].
	\end{aligned}
	\]
	
	For $z\in S_\rho$, the strip inclusion used in
	Proposition~\ref{prop:composition-strip} gives
	\[
	z+u_j(z)\in S_\tau.
	\]
	Indeed,
	\[
	|\Im(z+u_j(z))|
	<
	\log\rho+\bar d+R
	<
	\log\tau.
	\]
	Therefore
	\[
	|r_1(z+u_1(z))-r_2(z+u_1(z))|
	\le
	\Delta_r.
	\]
	Furthermore,
	\[
	|e^{i(z+u_j(z))}|
	=
	e^{-\Im(z+u_j(z))}
	\le
	\rho e^d.
	\]
	Thus the first term is bounded by
	\[
	\rho e^d\Delta_r.
	\]
	For the second term, we have
	\[r_2(z+u_1(z))-r_2(z+u_2(z))=
		\int_0^1 r_2'\bigl(z+u_2(z)+s(u_1(z)-u_2(z))\bigr)
		\bigl(u_1(z)-u_2(z)\bigr)\,ds .
	\]
	Since the intermediate arguments remain in $S_\tau$ and
	$\|r_2'\|_\tau\le K_\delta$, we obtain
	\[
	|r_2(z+u_1(z))-r_2(z+u_2(z))|
	\le
	K_\delta |u_1(z)-u_2(z)|.
	\]
	Moreover,
	\[
	|u_1(z)-u_2(z)|
	\le
	\|u_1-u_2\|_\rho
	\le
	C_u\Delta_r.
	\]
	Hence the second term is bounded by
	\[
	\rho e^d K_\delta C_u\Delta_r.
	\]
	Similarly,
	\[
	\begin{aligned}
		&e^{i(z+u_1(z))}-e^{i(z+u_2(z))}
		\\
		&\qquad =
		i\int_0^1
		e^{i(z+u_2(z)+s(u_1(z)-u_2(z)))}
		\bigl(u_1(z)-u_2(z)\bigr)\,ds .
	\end{aligned}
	\]
	The intermediate values satisfy the same bound on their imaginary
	parts, and therefore
	\[
	\left|
	e^{i(z+u_1(z))}-e^{i(z+u_2(z))}
	\right|
	\le
	\rho e^d C_u\Delta_r.
	\]
	Since
	\[
	|r_2(z+u_2(z))|
	\le
	\|r_2\|_\tau
	\le
	R_\delta,
	\]
	the third term is bounded by
	\[
	\rho e^d R_\delta C_u\Delta_r.
	\]
	
	Combining the three estimates gives
	\begin{equation}\label{eq:F-strip-bound}
		\sup_{z\in S_\rho}
		|F_1(z)-F_2(z)|
		\le
		M,
	\end{equation}
	where
	\[
	M
	:=
	\rho e^d
	\left[
	1+(K_\delta+R_\delta)C_u
	\right]\Delta_r.
	\]
	
	Write
	\[
	f_{r_1}(\zeta)-f_{r_2}(\zeta)
	=
	\sum_{k\ge1}c_k\zeta^k.
	\]
	Fix $1<\rho'<\rho$.  If $|\zeta|=\rho'$, write
	\[
	\zeta=e^{iz},
	\qquad
	\Im z=-\log\rho'.
	\]
	Then $z\in S_\rho$, and \eqref{eq:F-f-identification} and
	\eqref{eq:F-strip-bound} give
	\[
	|f_{r_1}(\zeta)-f_{r_2}(\zeta)|
	\le M.
	\]
	Cauchy's estimate therefore yields
	\[
	|c_k|\le M(\rho')^{-k}.
	\]
	Letting $\rho'\uparrow\rho$, we obtain
	\[
	|c_k|\le M\rho^{-k}.
	\]
	Consequently,
	\[
	\begin{aligned}
		\|f_{r_1}-f_{r_2}\|_{\rho_f}
		&=
		\sum_{k\ge1}|c_k|\rho_f^k
		\\
		&\le
		M\sum_{k\ge1}
		\left(\frac{\rho_f}{\rho}\right)^k
		\\
		&=
		M\frac{\rho_f}{\rho-\rho_f}.
	\end{aligned}
	\]
	Substituting the value of $M$ proves
	\eqref{eq:f-lipschitz} with $C_f$ as defined in
	\eqref{eq:Cf}.
\end{proof}

\end{document}
